\documentclass[11pt]{article}

\usepackage[a4paper,margin=1in]{geometry}
\usepackage{amsmath,amssymb,amsthm,mathtools}
\usepackage{enumitem}
\usepackage[hidelinks]{hyperref}

\newtheorem{theorem}{Theorem}[section]
\newtheorem{proposition}[theorem]{Proposition}
\newtheorem{lemma}[theorem]{Lemma}
\newtheorem{corollary}[theorem]{Corollary}
\theoremstyle{definition}
\newtheorem{definition}[theorem]{Definition}
\newtheorem{remark}[theorem]{Remark}

\newcommand{\R}{\mathbb R}
\newcommand{\Om}{\Omega}
\newcommand{\Fr}{\mathcal A}
\newcommand{\HH}{\mathcal H}
\newcommand{\dT}{\delta_T}
\newcommand{\rstar}{\rho_*}
\newcommand{\rdel}{\rho_\delta}
\newcommand{\wt}{w}
\newcommand{\eps}{\varepsilon}
\newcommand{\abs}[1]{\left|#1\right|}

\title{Raw Moving-Ball First Variation and Closure of the Bad-Density Step\\
       \large Plan 2 / Level-Set Rearrangement Route}
\author{}
\date{May 2026}

\begin{document}
\maketitle

\begin{abstract}
This note proves the remaining local theorem left by
\texttt{faber\_krahn\_remaining\_steps.tex}, with one correction: the right
moving-ball vector field is not the singular unit field \(e_z\), but the
affine homothetic field \((x-z)/\rho+a\).  It maps \(B_\rho(z)\) exactly to
\(B_{\rho+h}(z+ha)\), is globally Lipschitz, and pairs with the existing
homothetic radial-trace estimate.  We prove a finite-perimeter one-sided
first variation for
\[
   q(\rho):=\inf_{z\in\R^n}|E_\rho\Delta B_\rho(z)|.
\]
The estimate controls only the positive derivative \(q'_+\), which is exactly
what the endpoint trace needs.  Together with nestedness, the bad-density
set is paid at \(O(\delta_T)\), and the bounded sharp Saint--Venant estimate
follows from the level-set identity, same-centre Fusco--Julin, and the
boundary-layer transfer.
\end{abstract}

\tableofcontents

\section{Setup}

Let \(\Omega\subset B_R\) be bounded, \(|\Omega|=\omega_n\), and let \(u\) be
the torsion potential.  Set
\[
   E_\rho:=\{u>t(\rho)\},\qquad |E_\rho|=\omega_n\rho^n,\qquad
   \wt(\rho):=-t'(\rho),\qquad d\mu=\wt(\rho)\,d\rho .
\]
For a.e. \(\rho\), \(E_\rho\) has finite perimeter, \(|\nabla u|>0\) on
\(\partial^*E_\rho\) up to \(\HH^{n-1}\)-null sets, and
\[
   V_\rho=\frac{\wt(\rho)}{|\nabla u|},\qquad
   \int_{\partial^*E_\rho}V_\rho\,d\HH^{n-1}=P(B_\rho).
\]
Write
\[
   H_{z,\rho}(x):=\frac{(x-z)\cdot\nu_\rho(x)}{\rho},\qquad
   \bar V_\rho:=\frac{P(B_\rho)}{P(E_\rho)}.
\]
The exact level-set identity supplies
\begin{equation}\label{eq:budget}
   \int_{\rstar}^{\rdel}(D_H(t(\rho))+D_I(t(\rho)))\,d\mu(\rho)
   \le C(n,\rstar)\delta_T(\Omega).
\end{equation}
Talenti's profile-gap estimate supplies
\begin{equation}\label{eq:badtau}
   |B_\tau|\le C(n,\rstar)\delta_T,\qquad
   B_\tau:=\{\rho:\wt(\rho)<\tau_0\},\quad
   \tau_0:=\frac{\rstar}{2n}.
\end{equation}
Finally set
\[
   G_I:=\{\rho:D_I(t(\rho))\le\eta_I\}.
\]
On \(G_I\), the same-centre Fusco--Julin package and the oriented radial
trace lemma give, at a same-centre field \(z_\rho^{\rm sc}\),
\begin{align}
   |E_\rho\Delta B_\rho(z_\rho^{\rm sc})|^2
   &\le C(n,R,\rstar)D_I(t(\rho)), \label{eq:q-DI}\\
   \left(\int_{\partial^*E_\rho}|H_{z_\rho,\rho}-\bar V_\rho|\,d\HH^{n-1}\right)^2
   &\le C(n,R,\rstar)D_I(t(\rho)), \label{eq:H-DI}
\end{align}
where \(z_\rho=z_\rho^{\rm sc}\) in \eqref{eq:H-DI}.  The first inequality
implies the same bound for \(q(\rho)^2\), since \(q\) is an infimum over
centres.

\section{A Coarea Differentiation Lemma}

The proof uses the following standard consequence of BV coarea and Lebesgue
differentiation.  It is the precise replacement for smooth normal-flow
coordinates.

\begin{lemma}[Variable-thickness coarea differentiation]\label{lem:varcoarea}
For a.e. regular level \(s\) of \(u\), and every bounded Borel function
\(g\ge0\),
\begin{equation}\label{eq:varcoarea}
   \lim_{h\downarrow0}\frac1h
   \left|\{x:0<u(x)-s<hg(x)\}\right|
   =
   \int_{\partial^*\{u>s\}}\frac{g(x)}{|\nabla u(x)|}\,d\HH^{n-1}(x),
\end{equation}
and the analogous identity holds for \(\{x:0<s-u(x)<hg(x)\}\).
\end{lemma}

\begin{proof}
For simple \(g=\sum_{j=1}^N c_j\mathbf1_{A_j}\), the left side is a finite
sum of terms
\[
   \frac1h\int_s^{s+hc_j}
      \int_{\partial^*\{u>\tau\}}
      \frac{\mathbf1_{A_j}(x)}{|\nabla u(x)|}\,d\HH^{n-1}(x)\,d\tau ,
\]
by the BV coarea formula.  At every common Lebesgue point of the finitely
many coarea densities
\[
   \tau\mapsto
   \int_{\partial^*\{u>\tau\}}
      \frac{\mathbf1_{A_j}}{|\nabla u|}\,d\HH^{n-1},
\]
this converges to the right side of \eqref{eq:varcoarea}.  Approximate a
bounded Borel \(g\) from below and above by simple functions and use the
monotonicity of the slabs.  A countable monotone-class argument gives the
identity for all bounded Borel \(g\) outside one null set of levels.  The
lower slab is identical.
\end{proof}

\section{Raw Moving-Ball First Variation}

\begin{lemma}[Affine shell estimate]\label{lem:shell}
Fix a regular level \(\rho\) satisfying the differentiability and coarea
Lebesgue-point properties above.  Let \(z,a\in\R^n\), and define
\[
   T_h(x):=z+ha+\left(1+\frac{h}{\rho}\right)(x-z).
\]
Then, for \(h\downarrow0\),
\begin{equation}\label{eq:shell}
   \limsup_{h\downarrow0}
   \frac{|E_{\rho+h}\Delta T_h(E_\rho)|}{h}
   \le
   \int_{\partial^*E_\rho}
      |V_\rho-H_{z,\rho}-a\cdot\nu_\rho|\,d\HH^{n-1}.
\end{equation}
\end{lemma}

\begin{proof}
The map \(T_h={\rm id}+hW\) has the Lipschitz velocity
\[
   W(x)=\frac{x-z}{\rho}+a.
\]
Since \(u\in W^{2,p}_{\rm loc}(\Omega)\) for \(p>n\), hence
\(C^{1,\alpha}_{\rm loc}\), a countable compact-exhaustion of
\(\{u=t(\rho)\}\cap\Omega\) gives, up to an \(\HH^{n-1}\)-null boundary part,
\[
   u(T_h^{-1}x)
   =
   u(x)-h\nabla u(x)\cdot W(x)+o(h)
\]
uniformly on each compact patch.  Also
\[
   t(\rho+h)=t(\rho)-h\wt(\rho)+o(h).
\]
Therefore the symmetric difference is contained, modulo an \(o(h)\)-thick
slab, in the set where \(u(x)-t(\rho)\) lies between \(0\) and
\[
   h\bigl(\wt(\rho)+\nabla u(x)\cdot W(x)\bigr)
\]
or between that number and \(0\).  Applying Lemma~\ref{lem:varcoarea} to
the positive and negative parts gives
\[
   \limsup_{h\downarrow0}
   \frac{|E_{\rho+h}\Delta T_h(E_\rho)|}{h}
   \le
   \int_{\partial^*E_\rho}
   \frac{|\wt(\rho)+\nabla u\cdot W|}{|\nabla u|}\,d\HH^{n-1}.
\]
On \(\partial^*E_\rho\), the outer normal is
\(\nu_\rho=-\nabla u/|\nabla u|\), so
\[
   \frac{\wt+\nabla u\cdot W}{|\nabla u|}
   =
   V_\rho-W\cdot\nu_\rho
   =
   V_\rho-H_{z,\rho}-a\cdot\nu_\rho .
\]
This proves \eqref{eq:shell}.
\end{proof}

\begin{theorem}[One-sided first variation of the moving-ball distance]
\label{thm:qplus}
For a.e. \(\rho\in[\rstar,\rdel]\), let
\[
   q(\rho):=\inf_z |E_\rho\Delta B_\rho(z)|.
\]
Then
\begin{equation}\label{eq:qplus}
   D^+q(\rho):=
   \limsup_{h\downarrow0}\frac{q(\rho+h)-q(\rho)}{h}
   \le
   \frac{n}{\rstar}q(\rho)
   +
   \inf_{a\in\R^n}\int_{\partial^*E_\rho}
      |V_\rho-H_{z_\rho,\rho}-a\cdot\nu_\rho|\,d\HH^{n-1},
\end{equation}
where \(z_\rho\) is any minimiser of \(z\mapsto |E_\rho\Delta B_\rho(z)|\).
\end{theorem}

\begin{proof}
The minimiser exists because \(z\mapsto |E_\rho\Delta B_\rho(z)|\) is
continuous and tends to \(2|E_\rho|\) as \(|z|\to\infty\).  Fix such
\(z_\rho\), fix \(a\in\R^n\), and use the affine map \(T_h\) from
Lemma~\ref{lem:shell}.  Since
\[
   T_h(B_\rho(z_\rho))=B_{\rho+h}(z_\rho+ha),
\]
we have
\[
\begin{aligned}
   q(\rho+h)
   &\le |E_{\rho+h}\Delta T_h(B_\rho(z_\rho))|\\
   &\le |E_{\rho+h}\Delta T_h(E_\rho)|
      +|T_h(E_\rho)\Delta T_h(B_\rho(z_\rho))|\\
   &= |E_{\rho+h}\Delta T_h(E_\rho)|
      +\left(1+\frac{h}{\rho}\right)^n q(\rho).
\end{aligned}
\]
Subtracting \(q(\rho)\), dividing by \(h\), taking the limsup, and using
\(\rho\ge\rstar\) and Lemma~\ref{lem:shell}, gives
\[
   D^+q(\rho)
   \le \frac{n}{\rho}q(\rho)
   +
   \int_{\partial^*E_\rho}|V_\rho-H_{z_\rho,\rho}-a\cdot\nu_\rho|\,d\HH^{n-1}.
\]
Take the infimum over \(a\).
\end{proof}

\begin{remark}[Why one-sided is enough]
A two-sided derivative estimate would require controlling possible jumps of
the optimal centre.  The endpoint trace only transports information forward
from a Markov radius to \(\rho_\delta\).  Thus the positive derivative
\((q')_+\), supplied by \eqref{eq:qplus}, is the exact quantity needed.
\end{remark}

\section{Integrated Positive Kinetic Bound}

\begin{lemma}[Transfer from the FJ centre to a distance-minimising centre]
\label{lem:center-transfer}
Let \(\rho\in G_I\).  Let \(z_\rho^{\rm sc}\) be a same-centre
Fusco--Julin centre satisfying \eqref{eq:q-DI}--\eqref{eq:H-DI}, and let
\(z_\rho^{\rm min}\) minimise \(z\mapsto |E_\rho\Delta B_\rho(z)|\).  Then
\begin{equation}\label{eq:H-DI-min}
   \left(
   \int_{\partial^*E_\rho}
      |H_{z_\rho^{\rm min},\rho}-\bar V_\rho|\,d\HH^{n-1}
   \right)^2
   \le C(n,R,\rstar)D_I(t(\rho)).
\end{equation}
\end{lemma}

\begin{proof}
By minimality,
\[
   |E_\rho\Delta B_\rho(z_\rho^{\rm min})|
   \le |E_\rho\Delta B_\rho(z_\rho^{\rm sc})|
   \le C D_I(t(\rho))^{1/2}.
\]
Hence
\[
   |B_\rho(z_\rho^{\rm min})\Delta B_\rho(z_\rho^{\rm sc})|
   \le
   |E_\rho\Delta B_\rho(z_\rho^{\rm min})|
   +|E_\rho\Delta B_\rho(z_\rho^{\rm sc})|
   \le C D_I^{1/2}.
\]
For two equal balls with centres separated by \(d\), one has
\[
   |B_\rho(z_1)\Delta B_\rho(z_2)|
   \ge c(n)\rho^{n-1}\min\{d,\rho\}.
\]
Choosing \(\eta_I\) small in the definition of \(G_I\), this gives
\[
   |z_\rho^{\rm min}-z_\rho^{\rm sc}|
   \le C(n,\rstar)D_I(t(\rho))^{1/2}.
\]
The homothetic trace is affine in the centre:
\[
   H_{z_\rho^{\rm min},\rho}-H_{z_\rho^{\rm sc},\rho}
   =
   -\frac{(z_\rho^{\rm min}-z_\rho^{\rm sc})\cdot\nu_\rho}{\rho}.
\]
On \(G_I\), \(P(E_\rho)\le C(n,\rstar,R)\).  Therefore
\[
   \int_{\partial^*E_\rho}
      |H_{z_\rho^{\rm min},\rho}-H_{z_\rho^{\rm sc},\rho}|\,d\HH^{n-1}
   \le C(n,R,\rstar)D_I(t(\rho))^{1/2}.
\]
Combine this with \eqref{eq:H-DI} and square.
\end{proof}

\begin{lemma}[Unconditional Lipschitz bound for \(q\)]\label{lem:qLip}
For \(0<\rho\le\sigma\le1\),
\[
   |q(\sigma)-q(\rho)|
   \le 2\omega_n(\sigma^n-\rho^n)
   \le 2n\omega_n|\sigma-\rho|.
\]
\end{lemma}

\begin{proof}
Use the same centre for the two balls and the nestedness
\(E_\rho\subset E_\sigma\), \(B_\rho(z)\subset B_\sigma(z)\):
\[
   q(\sigma)\le q(\rho)+|E_\sigma\setminus E_\rho|+|B_\sigma\setminus B_\rho|
   =q(\rho)+2\omega_n(\sigma^n-\rho^n),
\]
and reverse the roles of \(\rho,\sigma\).
\end{proof}

\begin{proposition}[Positive kinetic estimate]\label{prop:positivekinetic}
Let \(J\subset[\rstar,\rdel]\) be an order-one interval.  Then
\[
   \int_J (q'(\rho)_+)^2\,d\rho
   \le C(n,R,\rstar,J)\,\delta_T(\Omega),
\]
where \(q'_+=\max\{q',0\}\), defined a.e. by the Lipschitz property.
\end{proposition}

\begin{proof}
Split \(J\) into
\[
   G:=J\cap G_I\cap G_\tau,\qquad
   \mathcal B:=J\setminus G .
\]
On \(G\), Theorem~\ref{thm:qplus}, \eqref{eq:q-DI}, and
\[
\begin{aligned}
   \inf_a\int|V_\rho-H_{z_\rho^{\rm min},\rho}-a\cdot\nu_\rho|
   &\le \int|V_\rho-\bar V_\rho|
        +\int|H_{z_\rho^{\rm min},\rho}-\bar V_\rho|
\end{aligned}
\]
give, after squaring,
\[
   (q'_+)^2
   \le C(n,R,\rstar)(D_H(t(\rho))+D_I(t(\rho))).
\]
Here
\[
   \left(\int|V_\rho-\bar V_\rho|\right)^2\le D_H(t(\rho))
\]
is the unconditional velocity-variance identity, and \eqref{eq:H-DI} controls
the homothetic term after the transfer Lemma~\ref{lem:center-transfer}.
Since \(G\subset G_\tau\), \(d\rho\le\tau_0^{-1}d\mu\),
so \eqref{eq:budget} gives
\[
   \int_G(q'_+)^2\,d\rho\le C\delta_T.
\]
On \(\mathcal B\), Lemma~\ref{lem:qLip} gives \(q'_+\le2n\omega_n\).  Also
\[
   |J\cap B_\tau|\le C\delta_T
\]
by \eqref{eq:badtau}, while
\[
   |J\cap G_\tau\cap G_I^c|
   \le \tau_0^{-1}\mu(G_I^c)
   \le \frac{C}{\eta_I}\int_JD_I\,d\mu
   \le C\delta_T.
\]
Thus \(|\mathcal B|\le C\delta_T\), and
\[
   \int_{\mathcal B}(q'_+)^2\,d\rho\le C\delta_T.
\]
\end{proof}

\section{Endpoint Trace and Bounded Stability}

\begin{lemma}[Moving-ball endpoint trace]\label{lem:endpoint}
Let \(J=[\rdel-L,\rdel]\subset[\rstar,\rdel]\), with \(L>0\) fixed.  Then
\[
   q(\rdel)^2\le C(n,R,\rstar,L)\,\delta_T(\Omega).
\]
\end{lemma}

\begin{proof}
First,
\[
   \int_{J\cap G}q(\rho)^2\,d\rho
   \le C\int_{J\cap G}D_I(t(\rho))\,d\rho
   \le C\int_JD_I(t(\rho))\,d\mu(\rho)
   \le C\delta_T.
\]
On \(J\setminus G\), the measure is \(O(\delta_T)\) by the proof of
Proposition~\ref{prop:positivekinetic}, and \(q\le2\omega_n\), so
\[
   \int_J q(\rho)^2\,d\rho\le C\delta_T.
\]
Hence there is \(\rho_0\in J\) with \(q(\rho_0)^2\le C L^{-1}\delta_T\).
Since \(\rho_0\le\rdel\), the one-sided transport gives
\[
   q(\rdel)
   \le q(\rho_0)+\int_{\rho_0}^{\rdel}q'_+(\rho)\,d\rho
   \le C\delta_T^{1/2}
      +L^{1/2}\left(\int_J(q'_+)^2\,d\rho\right)^{1/2}
   \le C\delta_T^{1/2},
\]
using Proposition~\ref{prop:positivekinetic}.
\end{proof}

\begin{theorem}[Bounded sharp Saint--Venant stability]\label{thm:boundedSV}
Assume the level-set identity, Talenti bad-density estimate, same-centre
Fusco--Julin package, and finite-perimeter oriented radial trace lemma in
the Plan~2 repository.  Then for bounded \(\Omega\subset B_R\),
\[
   \Fr(\Omega)^2\le C(n,R)\,\delta_T(\Omega).
\]
\end{theorem}

\begin{proof}
Choose
\[
   \rdel=1-\kappa\sqrt{\delta_T(\Omega)}
\]
with the usual fixed \(\kappa=\kappa(n,R)\) and \(\delta_T\) small enough
that the order-one window \(J=[\rdel-L,\rdel]\) lies in \([\rstar,\rdel]\).
Lemma~\ref{lem:endpoint} gives
\[
   \Fr(E_{\rdel})^2\le C(n,\rstar)q(\rdel)^2\le C\delta_T.
\]
Since \(E_{\rdel}\subset\Omega\),
\[
   |\Omega\setminus E_{\rdel}|=\omega_n(1-\rdel^n)
   \le C(n)\sqrt{\delta_T}.
\]
The elementary boundary-layer transfer gives
\[
   \Fr(\Omega)\le \Fr(E_{\rdel})+\frac{2}{\omega_n}|\Omega\setminus E_{\rdel}|
   \le C\sqrt{\delta_T}.
\]
Squaring proves the theorem.  The large-deficit case is absorbed by reducing
the constant, since \(\Fr(\Omega)\le2\).
\end{proof}

\section{What Changed}

The remaining step is closed in the bounded route, but with two corrections
to the earlier formulation.

\begin{enumerate}[leftmargin=2em]
\item The deformation field is homothetic:
\[
   W(x)=\frac{x-z}{\rho}+a,
\]
not \(e_z+a\).  This avoids all singular-centre issues and uses exactly the
homothetic trace already proved in Plan~2.
\item The first variation is one-sided:
\[
   D^+q\le Cq+\hbox{homothetic boundary action}.
\]
This is sufficient because the trace transports from a Markov radius forward
to \(\rho_\delta\).  A two-sided estimate is unnecessary.
\item Bad levels are harmless for \(q\), because nestedness gives a genuine
perimeter-free Lipschitz bound for the unscaled moving-ball distance.  The
invalid old claim concerned the rescaled sets \(\rho^{-1}E_\rho\), not \(q\).
\end{enumerate}

All constants are computable from the constants in the level-set identity,
Talenti profile gap, same-centre Fusco--Julin, the radial trace lemma, and
the explicit Lipschitz constants above.

\begin{thebibliography}{9}

\bibitem{AFP2000}
L.~Ambrosio, N.~Fusco, and D.~Pallara,
\emph{Functions of Bounded Variation and Free Discontinuity Problems},
Oxford University Press, 2000.

\bibitem{Maggi2012}
F.~Maggi,
\emph{Sets of Finite Perimeter and Geometric Variational Problems},
Cambridge University Press, 2012.

\bibitem{FJ2014}
N.~Fusco and V.~Julin,
\emph{A strong form of the quantitative isoperimetric inequality},
Calc. Var. Partial Differential Equations \textbf{50} (2014), 925--937.

\end{thebibliography}

\end{document}
